\chapter{Discussion}
\label{ch:discussion}

\section{Answering the research question}

\begin{quote}
\emph{How do Q-learning and SARSA compare when a warehouse robot must learn the shortest
collision-free route to a delivery point?}
\end{quote}

\subsection{They are not better or worse; they answer different questions}

Q-learning produced a deployed route of 19.40 steps against an optimum of 19. SARSA produced
33.34. On deployed quality, Q-learning wins decisively, and the difference is significant at
$p = 6.7\times10^{-7}$.

Q-learning accumulated 21{,}626 walkway incidents while learning; SARSA accumulated 5{,}068. On
the cost of learning, SARSA wins by a factor of 4.3, significant at $p = 1.9\times10^{-9}$.

Both statements are true simultaneously, and neither is a ranking of the algorithms. They are
answers to two different questions that are frequently conflated: \emph{how good is the policy you
end up with} and \emph{what did it cost to get there}.

\subsection{The mechanism, confirmed directly}

Chapter~\ref{ch:theory} derived the mechanism: Q-learning's target maximises over next actions, so
a hazard reachable only by a non-greedy action never enters the state value, while SARSA's target
carries an $\varepsilon/|\mathcal{A}|$-weighted share of it.

The sensitivity sweep confirms this without relying on inference. Q-learning's incident count is
numerically identical --- $6{,}522$ --- at hazard penalties of $-10$, $-25$, $-50$, $-100$ and
$-200$, and its route keeps a constant distance of 1.74 from the walkway. Scaling the cost of an
incident by a factor of twenty changes its behaviour by nothing whatsoever, because the penalty is
structurally invisible to its state values. SARSA's clearance, over the same sweep, rises
monotonically from 1.91 to 5.39.

It is rare for a derivation and a measurement to line up this cleanly, and it is worth stating that
the agreement is what gives confidence that the implementation is measuring the intended quantity
rather than an artefact.

\subsection{The effect is a response to risk, not a property of the algorithms}

Chapter~\ref{ch:dataset} reported a negative control that deserves more prominence than its
placement there suggests. On a floor built from scattered racking with no localised hazard, both
algorithms converged on the \emph{identical} 39-step optimum. There was no difference to measure.

The same conclusion appears at the other end of the sensitivity sweep: at zero hazard penalty the
clearance gap between the algorithms nearly closes.

The on-policy/off-policy distinction is therefore not a general characteristic that one algorithm
is safe and the other is greedy. It is a response to a specific structural property of the task ---
a high-penalty region adjacent to the optimal path. Where that property is absent, so is the
effect. Secondary accounts of Cliff Walking that present SARSA as ``the safe algorithm'' omit this
condition, and the omission matters: on a task without an adjacent hazard, choosing SARSA buys
nothing and costs route length.

\subsection{Which to use, and when}

The decision follows from where learning happens.

\keypoint{Learning in simulation, deploying the result.} Q-learning. Training incidents are free
because they occur in a simulator, and its asymptotic optimality is exactly what is wanted. This
describes most industrial practice.

\keypoint{Learning on a live floor among people.} SARSA, or an explicit safe-exploration method.
Every incident is real, and a 4.3-fold reduction in incidents during training is worth a 14-step
detour that can be re-optimised later.

\keypoint{Either, followed by $\varepsilon$ decay to zero.} Both converge to the optimal policy in
the limit under GLIE conditions \cite{singh2000}. The residual $\varepsilon = 0.05$ used here was
retained deliberately to keep the distinction visible.

\section{An uncomfortable comparison}

The heuristic controller recovered from the dataset --- move toward the goal --- achieved the
optimal 19 steps with zero incidents, equalling Q-learning and beating SARSA.

This should not be waved away. On this floor, that controller is optimal, because the main aisle
runs straight from the start to the goal and ``move toward the goal'' is therefore sufficient. The
learned agents spent 3{,}000 episodes to match what a four-line rule achieves immediately.

The honest reading is that the warehouse floor used here is a good instrument for isolating the
on-policy/off-policy distinction, and a poor advertisement for reinforcement learning as a
solution to this particular routing problem. A greedy heuristic fails when the direct direction is
blocked and a detour is required, which is precisely the situation the generality check across
twelve floors begins to probe --- but that check was not designed to compare against the heuristic,
and so this report does not claim the learned policies are superior to it. Establishing that would
require floors where the heuristic demonstrably fails, and is identified as future work.

\section{Threats to validity}

\subsection{Internal validity}

\keypoint{The reward constants are assumed.} The dataset supplied no usable reward signal, so the
step cost, bump penalty, hazard penalty and goal reward were chosen. The sensitivity sweep
quantifies the dependence on the one that drives the result, and shows the conclusion holds across
$-10$ to $-200$ while disappearing at $0$.

\keypoint{A constant step size.} $\alpha = 0.10$ violates the Robbins--Monro square-summability
condition, so neither algorithm converges to a point. Both use the same value, so the comparison is
unaffected, but asymptotic claims are not supported by this design.

\keypoint{Shared hyperparameters.} $\alpha$, $\gamma$ and the $\varepsilon$ schedule were fixed and
shared rather than tuned per algorithm. This is the correct choice for a controlled comparison, but
a schedule better suited to one algorithm could shift the balance.

\subsection{External validity}

\keypoint{The hazard is a construction.} The racking is real; the cross-aisle and the walkway are
not. Without them there is no effect, so this is not a neutral design decision and is not presented
as one. What the study establishes is the behaviour of the two algorithms \emph{given} such a
structure, together with evidence about what happens when it is absent.

\keypoint{One floor geometry.} All floors share the same three-row cross-aisle arrangement; only
the racking varies. A different hazard topology --- scattered restricted zones, or a walkway
crossing the route rather than running parallel to it --- might behave differently.

\keypoint{Tabular only.} 400 states are enumerable. Under function approximation, off-policy
bootstrapping interacts with the deadly triad \cite{sutton2018} in ways this study cannot speak to.

\keypoint{Deterministic dynamics.} Slip probability was zero. Stochastic transitions would be
expected to widen the gap, since SARSA's conservatism also hedges against randomness, but this was
not measured.

\subsection{Construct validity}

\keypoint{``Incidents'' are a proxy.} Counting walkway entries treats every incident as equally
costly, which a real safety analysis would not.

\keypoint{Online return mixes two things.} It aggregates route efficiency and incident cost into
one number. The report therefore also presents incident counts separately, since they are the
quantity an operator would actually care about.
